\begin{pastexam}[label={px:ncku-108-comm-ds-08}]{成大電通所 108 學年度 問答題第 5 題}
（問答；10 分）

You are given an empty hash table of size 7 and bucket capacity 1 that uses open addressing. The
following sequence of keys is to be inserted: (15, 17, 8, 23, 3, 5). Insert these keys using each of the
following hash functions and approaches. Show the resulting hash table step by step. If overflow occurs,
say so, indicate the element that causes the overflow. (10 points)
\begin{enumerate}[label=(\Alph*)]
\item h(x) = x \% 7; linear probing.
\item h(x) = x \% 7; double hashing with h2(x) = x / 7 + 1.
\end{enumerate}
\end{pastexam}

\begin{pastsoln}
表有 $7$ 格（編號 $0$ 到 $6$），每格放一個鍵。 \texttt{x / 7} 是整數除法。 兩種方法都沒有發生
overflow：$6$ 個鍵都放得進去（表沒滿，而且 (B) 的探測序列都能走到空格）。

\textbf{(A) 線性探測}：依序試 $h(x), h(x) + 1, h(x) + 2, \ldots \pmod 7$（資結書第~12.3~節）。
\begin{center}
\small
\begin{tabular}{cccl}
\toprule
鍵 & $h(x)$ & 試過的格子 & 放入 \\
\midrule
15 & 1 & 1 & [1] \\
17 & 3 & 3 & [3] \\
8 & 1 & 1（15）、2 & [2] \\
23 & 2 & 2（8）、3（17）、4 & [4] \\
3 & 3 & 3（17）、4（23）、5 & [5] \\
5 & 5 & 5（3）、6 & [6] \\
\bottomrule
\end{tabular}
\end{center}
\begin{center}
\begin{tabular}{|*{7}{c|}}
\hline
[0] & [1] & [2] & [3] & [4] & [5] & [6] \\ \hline
 & 15 & 8 & 17 & 23 & 3 & 5 \\ \hline
\end{tabular}
\end{center}
可以看到 $1$ 到 $6$ 格連成一串（主要群聚），後面插入的 23、3、5 都要多試好幾格。

\textbf{(B) 雙重雜湊}：第 $i$ 次試 $(h(x) + i \cdot h_2(x)) \bmod 7$，$h_2(x) = \lfloor x/7 \rfloor + 1$。
$7$ 是質數，$h_2$ 的值都在 $1$ 到 $6$ 之間，和 $7$ 互質，探測序列會走遍所有格子（資結書定理~12.4）。
\begin{center}
\small
\begin{tabular}{ccccl}
\toprule
鍵 & $h(x)$ & $h_2(x)$ & 試過的格子 & 放入 \\
\midrule
15 & 1 & 3 & 1 & [1] \\
17 & 3 & 3 & 3 & [3] \\
8 & 1 & 2 & 1（15）、$1 + 2 = 3$（17）、$1 + 4 = 5$ & [5] \\
23 & 2 & 4 & 2 & [2] \\
3 & 3 & 1 & 3（17）、4 & [4] \\
5 & 5 & 1 & 5（8）、6 & [6] \\
\bottomrule
\end{tabular}
\end{center}
\begin{center}
\begin{tabular}{|*{7}{c|}}
\hline
[0] & [1] & [2] & [3] & [4] & [5] & [6] \\ \hline
 & 15 & 23 & 17 & 3 & 8 & 5 \\ \hline
\end{tabular}
\end{center}
\end{pastsoln}
